\documentclass[12pt,a4paper]{article}

\usepackage[T2A]{fontenc}
\usepackage[utf8]{inputenc}
\usepackage[russian]{babel}
\usepackage{amsmath,amssymb}
\usepackage{enumitem}
\usepackage{array}
\usepackage{booktabs}
\usepackage{longtable}
\usepackage{geometry}
\usepackage{xcolor}
\usepackage{hyperref}
\usepackage{tikz}
\usetikzlibrary{arrows.meta,positioning,shapes.geometric}

\geometry{left=2.5cm,right=2cm,top=2cm,bottom=2cm}

\title{Билет №59 \\ \small Гибкие методологии. Agile-манифест, Scrum Guide. История и основные положения экстремального программирования. Современные гибкие методологии. (СпБГУ)}
\author{Сериков Владислав Александрович}
\date{16 мая 2026}

\begin{document}

\maketitle
\tableofcontents
\newpage

\section{Предпосылки появления гибких методологий}

\textbf{Гибкие методологии разработки программного обеспечения} --- семейство подходов к управлению разработкой, основанных на итеративной и инкрементальной поставке продукта, тесном взаимодействии с заказчиком, адаптации к изменениям и самоорганизации команды.

Гибкие подходы возникли как реакция на ограничения классических каскадных моделей. В каскадной модели предполагается последовательное прохождение стадий:
\[
\text{требования} \rightarrow \text{проектирование} \rightarrow \text{реализация} \rightarrow \text{тестирование} \rightarrow \text{внедрение}.
\]
Такая схема удобна при стабильных требованиях, но плохо работает в условиях неопределённости, когда:
\begin{itemize}
    \item заказчик не может заранее полностью сформулировать требования;
    \item рынок и технологии быстро меняются;
    \item ценность продукта можно проверить только после его использования;
    \item позднее обнаружение ошибок стоит дорого.
\end{itemize}

Гибкие методологии используют другой принцип:
\[
\text{короткий цикл} \rightarrow \text{работающий инкремент} \rightarrow \text{обратная связь} \rightarrow \text{адаптация}.
\]

\subsection{Базовые понятия}

\begin{description}
    \item[Итерация] ограниченный по времени цикл работы, в течение которого команда выполняет анализ, проектирование, реализацию, тестирование и демонстрацию результата.
    \item[Инкремент] добавленная к продукту часть функциональности, пригодная к проверке и потенциальному выпуску.
    \item[Бэклог] упорядоченный список работ, требований, пользовательских историй, дефектов и технических задач.
    \item[Пользовательская история] краткое описание потребности пользователя в формате:
    \[
    \text{Как } \langle \text{роль} \rangle
    \text{ я хочу } \langle \text{действие} \rangle,
    \text{ чтобы } \langle \text{ценность} \rangle.
    \]
    \item[Definition of Done] явные критерии готовности элемента продукта.
    \item[Самоорганизующаяся команда] команда, самостоятельно выбирающая способ достижения поставленной цели.
\end{description}

\section{Agile-манифест}

\subsection{История Agile-манифеста}

Agile-манифест был сформулирован в 2001 году группой специалистов по разработке ПО, которые представляли разные лёгковесные методологии: Extreme Programming, Scrum, Crystal, Adaptive Software Development, Feature-Driven Development и другие.

Целью Agile-манифеста было не создание единственного процесса, а фиксация общих ценностей, которые отличали гибкие подходы от тяжеловесных, бюрократических и чрезмерно плановых моделей.

\subsection{Четыре ценности Agile}

Agile-манифест формулирует четыре ключевые ценности:

\begin{enumerate}
    \item \textbf{Люди и взаимодействие важнее процессов и инструментов.}

    Процессы и инструменты полезны, но они не должны заменять коммуникацию, ответственность и профессиональное суждение команды.

    \item \textbf{Работающий продукт важнее исчерпывающей документации.}

    Документация нужна, но главным доказательством прогресса является работающая система.

    \item \textbf{Сотрудничество с заказчиком важнее согласования условий контракта.}

    В гибкой разработке заказчик рассматривается не как внешний контролёр, а как постоянный участник уточнения продукта.

    \item \textbf{Готовность к изменениям важнее следования первоначальному плану.}

    Планирование необходимо, но план не должен препятствовать созданию ценности при изменении условий.
\end{enumerate}

Важно понимать, что Agile не отрицает процессы, инструменты, документацию, контракты и планы. Манифест утверждает, что элементы слева имеют больший приоритет, чем элементы справа.

\subsection{Двенадцать принципов Agile}

\begin{enumerate}
    \item Наивысший приоритет --- удовлетворение заказчика за счёт ранней и непрерывной поставки ценного программного обеспечения.
    \item Изменения требований приветствуются даже на поздних стадиях разработки.
    \item Работающий продукт следует выпускать часто: от нескольких недель до нескольких месяцев, предпочтительно с коротким интервалом.
    \item Представители бизнеса и разработчики должны ежедневно работать вместе.
    \item Проекты строятся вокруг мотивированных людей; им дают среду, поддержку и доверие.
    \item Наиболее эффективный способ передачи информации --- личное общение.
    \item Работающий продукт --- основной показатель прогресса.
    \item Гибкие процессы способствуют устойчивому темпу разработки.
    \item Постоянное внимание к техническому совершенству и качественному проектированию повышает гибкость.
    \item Простота --- искусство минимизации лишней работы --- крайне важна.
    \item Лучшие архитектуры, требования и проектные решения рождаются в самоорганизующихся командах.
    \item Команда регулярно анализирует способы повышения эффективности и корректирует своё поведение.
\end{enumerate}

\subsection{Agile как зонтичное понятие}

Agile --- не конкретная методология, а совокупность ценностей и принципов. Конкретными реализациями Agile являются, например:
\begin{itemize}
    \item Scrum;
    \item Extreme Programming;
    \item Kanban;
    \item Lean Software Development;
    \item Crystal;
    \item Feature-Driven Development;
    \item Dynamic Systems Development Method;
    \item Disciplined Agile;
    \item SAFe, LeSS, Nexus как подходы к масштабированию.
\end{itemize}

\section{Scrum}

\subsection{Общая характеристика Scrum}

\textbf{Scrum} --- лёгковесный фреймворк для разработки, поставки и сопровождения сложных продуктов. Scrum не предписывает детальных инженерных практик, а задаёт минимальную структуру ролей, событий, артефактов и правил эмпирического управления.

Scrum основан на эмпиризме и бережливом мышлении.

\subsection{Эмпиризм в Scrum}

Эмпирическое управление предполагает, что знания возникают из опыта, а решения принимаются на основе наблюдаемых результатов.

Три опоры эмпиризма:

\begin{enumerate}
    \item \textbf{Прозрачность.}

    Существенные аспекты процесса и продукта должны быть видимы участникам. Например, состояние бэклога, цель спринта, критерии готовности и прогресс команды должны быть понятны.

    \item \textbf{Инспекция.}

    Команда регулярно проверяет артефакты и ход работы, чтобы выявлять отклонения.

    \item \textbf{Адаптация.}

    Если обнаружено отклонение или новая информация, процесс и продукт корректируются.
\end{enumerate}

\subsection{Ценности Scrum}

Scrum выделяет пять ценностей:

\begin{enumerate}
    \item \textbf{Commitment} --- приверженность достижению целей.
    \item \textbf{Focus} --- фокусировка на работе спринта и целях команды.
    \item \textbf{Openness} --- открытость относительно работы и проблем.
    \item \textbf{Respect} --- уважение к людям и их компетенциям.
    \item \textbf{Courage} --- смелость делать правильные вещи и обсуждать сложные вопросы.
\end{enumerate}

\subsection{Команда Scrum}

Scrum-команда состоит из:
\begin{itemize}
    \item Product Owner;
    \item Scrum Master;
    \item Developers.
\end{itemize}

Scrum-команда является небольшой, кросс-функциональной и самоорганизующейся. В Scrum Guide 2020 подчёркивается идея одной команды, сфокусированной на одном продукте.

\subsubsection{Product Owner}

\textbf{Product Owner} отвечает за максимизацию ценности продукта, создаваемого командой. Основные обязанности:

\begin{itemize}
    \item формирование и поддержание Product Backlog;
    \item упорядочивание элементов бэклога;
    \item обеспечение ясности Product Goal;
    \item взаимодействие со стейкхолдерами;
    \item принятие решений о приоритетах.
\end{itemize}

Product Owner --- единственный человек, отвечающий за управление Product Backlog. Это не комитет, хотя Product Owner может учитывать мнения многих заинтересованных сторон.

\subsubsection{Scrum Master}

\textbf{Scrum Master} отвечает за понимание и применение Scrum. Он служит команде, Product Owner и организации.

Основные функции Scrum Master:
\begin{itemize}
    \item помощь команде в самоорганизации;
    \item устранение препятствий;
    \item фасилитация событий Scrum;
    \item обучение команды и организации принципам Scrum;
    \item защита команды от неэффективных вмешательств;
    \item помощь Product Owner в управлении бэклогом.
\end{itemize}

Scrum Master не является классическим начальником команды и не распределяет задачи директивно.

\subsubsection{Developers}

\textbf{Developers} --- участники Scrum-команды, создающие каждый спринт пригодный к использованию инкремент.

Они отвечают за:
\begin{itemize}
    \item планирование работы спринта;
    \item создание Sprint Backlog;
    \item обеспечение качества согласно Definition of Done;
    \item ежедневную адаптацию плана;
    \item профессиональное выполнение инженерной работы.
\end{itemize}

\subsection{События Scrum}

Scrum использует фиксированный набор событий, каждое из которых служит инспекции и адаптации.

\subsubsection{Sprint}

\textbf{Sprint} --- контейнер для всех остальных событий Scrum. Это ограниченный по времени период, обычно не более одного месяца, в течение которого создаётся инкремент продукта.

Спринт включает:
\[
\text{Sprint Planning} \rightarrow
\text{Daily Scrum} \rightarrow
\text{работа над инкрементом} \rightarrow
\text{Sprint Review} \rightarrow
\text{Sprint Retrospective}.
\]

\subsubsection{Sprint Planning}

На планировании спринта команда отвечает на три вопроса:
\begin{enumerate}
    \item Почему этот спринт ценен?
    \item Что может быть сделано в этом спринте?
    \item Как выбранная работа будет выполнена?
\end{enumerate}

Результатом является Sprint Goal и Sprint Backlog.

\subsubsection{Daily Scrum}

\textbf{Daily Scrum} --- ежедневное событие для Developers продолжительностью обычно до 15 минут. Его цель --- инспекция прогресса к Sprint Goal и адаптация Sprint Backlog.

Daily Scrum не является отчётом начальнику. Это рабочая встреча команды для синхронизации.

\subsubsection{Sprint Review}

\textbf{Sprint Review} проводится в конце спринта для инспекции результата и адаптации Product Backlog. На встрече команда демонстрирует инкремент, обсуждает изменения рынка, требований, бюджета, сроков и дальнейшие шаги.

\subsubsection{Sprint Retrospective}

\textbf{Sprint Retrospective} --- встреча, посвящённая улучшению процесса работы. Команда обсуждает:
\begin{itemize}
    \item что прошло хорошо;
    \item какие проблемы возникли;
    \item что можно улучшить;
    \item какие конкретные действия будут предприняты в следующем спринте.
\end{itemize}

\subsection{Артефакты Scrum}

\subsubsection{Product Backlog}

\textbf{Product Backlog} --- упорядоченный список всего, что может потребоваться продукту. Он является единственным источником работы для Scrum-команды.

Product Backlog содержит:
\begin{itemize}
    \item пользовательские истории;
    \item дефекты;
    \item технические задачи;
    \item исследования;
    \item улучшения.
\end{itemize}

Связанное обязательство --- \textbf{Product Goal}, то есть долгосрочная цель продукта.

\subsubsection{Sprint Backlog}

\textbf{Sprint Backlog} --- план Developers на спринт. Он включает:
\begin{itemize}
    \item Sprint Goal;
    \item выбранные элементы Product Backlog;
    \item план их реализации.
\end{itemize}

Связанное обязательство --- \textbf{Sprint Goal}.

\subsubsection{Increment}

\textbf{Increment} --- конкретный шаг к Product Goal. Инкремент должен соответствовать Definition of Done и быть пригодным к использованию.

Связанное обязательство --- \textbf{Definition of Done}.

\subsection{Алгоритм Scrum-цикла}

\subsubsection*{Шаги}

\begin{enumerate}
    \item Product Owner поддерживает Product Backlog и упорядочивает элементы по ценности, рискам и зависимостям.
    \item На Sprint Planning команда выбирает цель спринта и элементы бэклога.
    \item Developers формируют Sprint Backlog и планируют реализацию.
    \item В течение спринта команда разрабатывает инкремент.
    \item Каждый день Developers проводят Daily Scrum и корректируют план.
    \item В конце спринта проводится Sprint Review: результат демонстрируется и обсуждается со стейкхолдерами.
    \item После этого проводится Sprint Retrospective: команда улучшает процесс.
    \item Product Backlog уточняется, и начинается следующий спринт.
\end{enumerate}

\subsubsection*{Схема}

\begin{center}
\begin{tikzpicture}[
    node distance=1.7cm,
    box/.style={rectangle, draw, rounded corners, align=center, minimum width=3.2cm, minimum height=1cm},
    arrow/.style={-{Latex}, thick}
]
\node[box] (pb) {Product\\Backlog};
\node[box, right=of pb] (plan) {Sprint\\Planning};
\node[box, right=of plan] (sprint) {Sprint\\1--4 недели};
\node[box, right=of sprint] (inc) {Increment};
\node[box, below=of sprint] (daily) {Daily\\Scrum};
\node[box, below=of inc] (review) {Sprint\\Review};
\node[box, below=of plan] (retro) {Sprint\\Retrospective};

\draw[arrow] (pb) -- (plan);
\draw[arrow] (plan) -- (sprint);
\draw[arrow] (sprint) -- (inc);
\draw[arrow] (sprint) -- (daily);
\draw[arrow] (daily) -- (sprint);
\draw[arrow] (inc) -- (review);
\draw[arrow] (review) -- (retro);
\draw[arrow] (retro) -- (pb);
\end{tikzpicture}
\end{center}

\subsubsection*{Минимальный пример}

Пусть разрабатывается интернет-магазин.

\begin{itemize}
    \item Product Goal: пользователь может выбрать товар, оплатить заказ и получить подтверждение.
    \item Элемент Product Backlog:
    \[
    \text{Как покупатель, я хочу добавить товар в корзину, чтобы оформить заказ позже.}
    \]
    \item Sprint Goal: реализовать базовую корзину.
    \item Sprint Backlog:
    \begin{enumerate}
        \item создать таблицу корзин;
        \item реализовать API добавления товара;
        \item реализовать API удаления товара;
        \item покрыть сценарии тестами;
        \item обновить интерфейс.
    \end{enumerate}
    \item Increment: пользователь может добавить товар в корзину, увидеть список товаров и удалить товар.
\end{itemize}

\section{Extreme Programming}

\subsection{Общая характеристика}

\textbf{Extreme Programming, XP} --- гибкая методология разработки программного обеспечения, ориентированная на высокое качество кода, частую обратную связь, тесное взаимодействие с заказчиком и технические практики, снижающие стоимость изменений.

Название ``экстремальное программирование'' связано с идеей доведения полезных инженерных практик до ``экстремального'' уровня:
\begin{itemize}
    \item если код-ревью полезно, то код проверяется постоянно через парное программирование;
    \item если тестирование полезно, то тесты пишутся до кода;
    \item если интеграция полезна, то интеграция выполняется часто;
    \item если простота полезна, то проектируется только то, что действительно нужно сейчас.
\end{itemize}

\subsection{История XP}

XP сформировалось в 1990-е годы, прежде всего в работах Кента Бека. Важным проектом для становления XP был Chrysler Comprehensive Compensation System, где применялись короткие итерации, тесная связь с пользователями, автоматизированные тесты и коллективная работа над кодом.

В 1999 году Кент Бек опубликовал книгу \textit{Extreme Programming Explained: Embrace Change}, в которой систематизировал ценности, принципы и практики XP.

XP оказало большое влияние на всё Agile-движение. Многие практики, которые сегодня воспринимаются как стандарт индустрии, были популяризованы именно XP:
\begin{itemize}
    \item test-driven development;
    \item continuous integration;
    \item pair programming;
    \item refactoring;
    \item user stories;
    \item small releases.
\end{itemize}

\subsection{Ценности XP}

Классические ценности XP:

\begin{enumerate}
    \item \textbf{Communication} --- постоянная коммуникация между разработчиками, заказчиком и пользователями.
    \item \textbf{Simplicity} --- реализация наиболее простого решения, достаточного для текущих требований.
    \item \textbf{Feedback} --- быстрая обратная связь от тестов, заказчика, пользователей и команды.
    \item \textbf{Courage} --- готовность менять код, удалять лишнее, признавать ошибки и принимать изменения.
    \item \textbf{Respect} --- уважение к участникам команды и их профессиональному вкладу.
\end{enumerate}

\subsection{Принципы XP}

К основным принципам XP относятся:
\begin{itemize}
    \item быстрая обратная связь;
    \item предположение простоты;
    \item постепенные изменения;
    \item принятие изменений;
    \item качественная работа;
    \item экономия усилий;
    \item взаимная выгода;
    \item самоорганизация;
    \item ответственность.
\end{itemize}

\subsection{Практики XP}

\subsubsection{Планирование}

\begin{description}
    \item[Planning Game] совместное планирование заказчиком и разработчиками. Заказчик определяет ценность функций, разработчики оценивают стоимость.
    \item[Small Releases] частые небольшие выпуски работающего продукта.
    \item[On-site Customer] постоянное присутствие представителя заказчика или доступность доменного эксперта.
\end{description}

\subsubsection{Проектирование}

\begin{description}
    \item[Simple Design] проектирование максимально простого решения.
    \item[System Metaphor] использование общей метафоры системы, помогающей команде понимать архитектуру.
    \item[Refactoring] регулярное улучшение внутренней структуры кода без изменения внешнего поведения.
\end{description}

\subsubsection{Кодирование}

\begin{description}
    \item[Pair Programming] два разработчика работают за одним рабочим местом: один пишет код, другой анализирует и предлагает улучшения.
    \item[Collective Code Ownership] любой разработчик может улучшать любой участок кода.
    \item[Coding Standards] единый стиль кодирования.
    \item[Continuous Integration] частая интеграция изменений в общую кодовую базу.
\end{description}

\subsubsection{Тестирование}

\begin{description}
    \item[Test-Driven Development] разработка через тестирование: сначала пишется тест, затем минимальный код, затем выполняется рефакторинг.
    \item[Acceptance Tests] приёмочные тесты, описывающие ожидаемое поведение с точки зрения заказчика.
\end{description}

\subsection{Алгоритм TDD}

\textbf{Test-Driven Development} --- практика, при которой разработка ведётся короткими циклами:
\[
\text{Red} \rightarrow \text{Green} \rightarrow \text{Refactor}.
\]

\subsubsection*{Шаги}

\begin{enumerate}
    \item Выбрать маленькое требование или поведение.
    \item Написать автоматический тест, который проверяет это поведение.
    \item Запустить тест и убедиться, что он падает. Это стадия \textbf{Red}.
    \item Написать минимальный код, достаточный для прохождения теста.
    \item Запустить тесты и убедиться, что они проходят. Это стадия \textbf{Green}.
    \item Улучшить структуру кода без изменения поведения. Это стадия \textbf{Refactor}.
    \item Повторить цикл для следующего поведения.
\end{enumerate}

\subsubsection*{Схема}

\begin{center}
\begin{tikzpicture}[
    node distance=2.5cm,
    state/.style={circle, draw, minimum size=2.2cm, align=center},
    arrow/.style={-{Latex}, thick}
]
\node[state, fill=red!20] (red) {Red\\тест падает};
\node[state, right=of red, fill=green!20] (green) {Green\\тест проходит};
\node[state, below=of green, fill=blue!15] (refactor) {Refactor\\улучшение};
\draw[arrow] (red) -- (green);
\draw[arrow] (green) -- (refactor);
\draw[arrow] (refactor) -- (red);
\end{tikzpicture}
\end{center}

\subsubsection*{Минимальный пример}

Требуется реализовать функцию сложения двух чисел.

\paragraph{Шаг 1. Тест.}

\begin{verbatim}
def test_add_two_numbers():
    assert add(2, 3) == 5
\end{verbatim}

Тест падает, потому что функция \texttt{add} ещё не реализована.

\paragraph{Шаг 2. Минимальная реализация.}

\begin{verbatim}
def add(a, b):
    return a + b
\end{verbatim}

Тест проходит.

\paragraph{Шаг 3. Рефакторинг.}

В данном примере структура уже проста, поэтому рефакторинг не требуется. На более сложных примерах можно переименовать переменные, выделить функции, убрать дублирование.

\subsection{Алгоритм парного программирования}

\subsubsection*{Шаги}

\begin{enumerate}
    \item Два разработчика выбирают задачу.
    \item Один становится \textbf{driver}: пишет код.
    \item Второй становится \textbf{navigator}: следит за направлением решения, качеством, тестами и возможными ошибками.
    \item Через короткий интервал роли меняются.
    \item Код считается общим результатом пары.
\end{enumerate}

\subsubsection*{Минимальный пример}

При реализации функции поиска скидки:
\begin{itemize}
    \item driver пишет тест на скидку для постоянного клиента;
    \item navigator замечает, что нужно также проверить граничный случай нулевой суммы заказа;
    \item после прохождения теста роли меняются;
    \item второй разработчик реализует обработку промокода, первый проверяет простоту решения.
\end{itemize}

\subsection{Сильные и слабые стороны XP}

\subsubsection*{Преимущества}

\begin{itemize}
    \item высокое качество кода;
    \item быстрая обратная связь;
    \item снижение стоимости изменений;
    \item хорошая пригодность для неопределённых требований;
    \item постоянное улучшение архитектуры через рефакторинг.
\end{itemize}

\subsubsection*{Ограничения}

\begin{itemize}
    \item требуется высокая дисциплина команды;
    \item заказчик должен быть доступен;
    \item парное программирование может встретить организационное сопротивление;
    \item без автоматизированных тестов XP теряет значительную часть эффективности;
    \item не все практики легко внедряются в больших распределённых организациях.
\end{itemize}

\section{Современные гибкие методологии и подходы}

\subsection{Kanban}

\textbf{Kanban} --- метод управления потоком работ, основанный на визуализации задач, ограничении незавершённой работы и постоянном улучшении.

Kanban не требует фиксированных итераций. Вместо этого используется непрерывный поток задач.

\subsubsection{Основные принципы Kanban}

\begin{enumerate}
    \item Визуализировать поток работы.
    \item Ограничивать количество незавершённой работы.
    \item Управлять потоком.
    \item Делать политики процесса явными.
    \item Вводить циклы обратной связи.
    \item Улучшать процесс эволюционно.
\end{enumerate}

\subsubsection{Kanban-доска}

Простейшая доска:

\[
\begin{array}{|c|c|c|}
\hline
\text{To Do} & \text{In Progress} & \text{Done} \\
\hline
\text{Задача A} & \text{Задача B} & \text{Задача C} \\
\hline
\end{array}
\]

Если для колонки \textit{In Progress} задан WIP-лимит 2, то в ней не может находиться более двух задач одновременно.

\subsubsection{Алгоритм работы Kanban}

\begin{enumerate}
    \item Все задачи помещаются в начальную колонку.
    \item Команда берёт новую задачу только при наличии свободного WIP-слота.
    \item Задача перемещается по колонкам согласно состоянию.
    \item Измеряются метрики потока: lead time, cycle time, throughput.
    \item Узкие места устраняются через изменение процесса.
\end{enumerate}

\subsection{Scrumban}

\textbf{Scrumban} --- гибрид Scrum и Kanban. Обычно он сочетает:
\begin{itemize}
    \item роли и планирование из Scrum;
    \item визуализацию потока и WIP-лимиты из Kanban;
    \item более гибкое отношение к длительности итераций.
\end{itemize}

Scrumban полезен для команд сопровождения, DevOps-команд и продуктовых команд, где часть работы плановая, а часть поступает нерегулярно.

\subsection{Lean Software Development}

\textbf{Lean Software Development} переносит идеи бережливого производства в разработку ПО.

Основные принципы:
\begin{enumerate}
    \item устранять потери;
    \item усиливать обучение;
    \item принимать решения как можно позже, но не слишком поздно;
    \item поставлять как можно быстрее;
    \item уважать людей;
    \item встраивать качество;
    \item оптимизировать целое, а не локальные части.
\end{enumerate}

Примеры потерь в разработке:
\begin{itemize}
    \item незавершённые задачи;
    \item лишняя функциональность;
    \item ожидание согласований;
    \item переключение контекста;
    \item дефекты;
    \item повторное выполнение работы;
    \item избыточная документация.
\end{itemize}

\subsection{Feature-Driven Development}

\textbf{Feature-Driven Development, FDD} --- методология, ориентированная на проектирование и реализацию функций.

Типичный процесс FDD:
\begin{enumerate}
    \item разработать общую модель предметной области;
    \item построить список функций;
    \item спланировать разработку по функциям;
    \item спроектировать функцию;
    \item реализовать функцию.
\end{enumerate}

Функция обычно формулируется как небольшое клиентски значимое действие:
\[
\text{вычислить итоговую стоимость заказа},
\quad
\text{зарегистрировать нового пользователя}.
\]

\subsection{Crystal}

\textbf{Crystal} --- семейство методологий, предложенное Алистером Кокберном. Основная идея: процесс должен зависеть от размера команды, критичности системы и организационного контекста.

Crystal делает акцент на:
\begin{itemize}
    \item людях и коммуникации;
    \item частой поставке;
    \item отражательном улучшении;
    \item минимально достаточном процессе.
\end{itemize}

\subsection{Dynamic Systems Development Method}

\textbf{DSDM} --- agile-подход, исторически связанный с быстрой разработкой приложений. Он ориентирован на фиксированные сроки и стоимость при варьируемом объёме работ.

В DSDM часто используется приоритизация MoSCoW:
\begin{itemize}
    \item \textbf{Must have} --- обязательно должно быть;
    \item \textbf{Should have} --- желательно;
    \item \textbf{Could have} --- можно сделать при наличии ресурсов;
    \item \textbf{Won't have now} --- не входит в текущий релиз.
\end{itemize}

\subsection{Disciplined Agile}

\textbf{Disciplined Agile} --- набор инструментов и практик, который помогает выбирать подходящий жизненный цикл и процесс под конкретный контекст. В отличие от Scrum, он менее prescriptive и предлагает варианты.

Disciplined Agile рассматривает разные жизненные циклы:
\begin{itemize}
    \item agile;
    \item lean;
    \item continuous delivery;
    \item exploratory;
    \item program-level lifecycle.
\end{itemize}

\subsection{SAFe}

\textbf{Scaled Agile Framework, SAFe} --- фреймворк масштабирования Agile для крупных организаций. Он предназначен для синхронизации множества команд, работающих над большими продуктами и портфелями.

Ключевые элементы SAFe:
\begin{itemize}
    \item Agile Release Train;
    \item Program Increment Planning;
    \item Lean Portfolio Management;
    \item роли уровня команды, программы и портфеля;
    \item синхронизация поставки на уровне нескольких команд.
\end{itemize}

SAFe полезен в крупных компаниях, но часто критикуется за сложность, формализацию и риск превращения Agile в бюрократический процесс.

\subsection{LeSS}

\textbf{Large-Scale Scrum, LeSS} --- подход масштабирования Scrum, который старается сохранить простоту Scrum при работе нескольких команд над одним продуктом.

Основные идеи LeSS:
\begin{itemize}
    \item один Product Backlog;
    \item один Product Owner;
    \item несколько команд;
    \item общий Sprint;
    \item общий потенциально поставляемый инкремент.
\end{itemize}

\subsection{Nexus}

\textbf{Nexus} --- фреймворк масштабирования Scrum для нескольких Scrum-команд, работающих над одним продуктом. Nexus вводит дополнительные события и артефакты для управления зависимостями и интеграцией.

\subsection{Agile и DevOps}

Современная гибкая разработка часто связана с DevOps.

\textbf{DevOps} --- культурный и инженерный подход, объединяющий разработку, эксплуатацию, автоматизацию поставки и наблюдаемость системы.

Связанные практики:
\begin{itemize}
    \item continuous integration;
    \item continuous delivery;
    \item infrastructure as code;
    \item автоматизированное тестирование;
    \item мониторинг и observability;
    \item быстрый откат изменений;
    \item feature flags.
\end{itemize}

Agile отвечает на вопрос, как организовать создание ценности, а DevOps --- как быстро и надёжно доставлять эту ценность пользователю.

\section{Сравнение гибких подходов}

\begin{longtable}{p{3cm}p{4cm}p{4cm}p{4cm}}
\toprule
\textbf{Подход} & \textbf{Основной фокус} & \textbf{Сильные стороны} & \textbf{Ограничения} \\
\midrule
Scrum & Управление продуктовой разработкой через спринты & Простая структура, регулярная обратная связь, ясные роли & Не задаёт инженерных практик \\
\midrule
XP & Инженерные практики и качество кода & TDD, CI, рефакторинг, парное программирование & Требует высокой дисциплины \\
\midrule
Kanban & Управление потоком задач & Гибкость, визуализация, WIP-лимиты & Меньше структуры для продуктового планирования \\
\midrule
Scrumban & Комбинация Scrum и Kanban & Подходит для смешанных потоков работ & Может стать неопределённым процессом без правил \\
\midrule
Lean & Устранение потерь и оптимизация целого & Хорош для организационных улучшений & Требует системного мышления \\
\midrule
SAFe & Масштабирование Agile в крупных организациях & Синхронизация многих команд & Сложность и риск бюрократизации \\
\midrule
LeSS & Масштабирование Scrum с минимальной сложностью & Сохраняет простоту Scrum & Требует организационных изменений \\
\bottomrule
\end{longtable}

\section{Типичные артефакты гибкой разработки}

\subsection{Пользовательские истории}

Пользовательская история описывает требование с точки зрения пользователя.

\[
\text{Как пользователь, я хочу восстановить пароль, чтобы вернуть доступ к аккаунту.}
\]

Хорошая история соответствует критериям INVEST:
\begin{itemize}
    \item \textbf{Independent} --- независимая;
    \item \textbf{Negotiable} --- обсуждаемая;
    \item \textbf{Valuable} --- ценная;
    \item \textbf{Estimable} --- оцениваемая;
    \item \textbf{Small} --- небольшая;
    \item \textbf{Testable} --- проверяемая.
\end{itemize}

\subsection{Критерии приёмки}

Критерии приёмки задают условия, при которых история считается реализованной.

Пример:

\begin{verbatim}
Дано: пользователь находится на странице входа
Когда: он нажимает "Забыли пароль?"
И: вводит зарегистрированный email
Тогда: система отправляет письмо со ссылкой восстановления
\end{verbatim}

\subsection{Оценивание задач}

В гибких командах часто используются относительные оценки:
\begin{itemize}
    \item story points;
    \item t-shirt sizes;
    \item planning poker.
\end{itemize}

Story points оценивают не часы, а относительную сложность, включая:
\begin{itemize}
    \item объём работы;
    \item неопределённость;
    \item риски;
    \item техническую сложность.
\end{itemize}

\section{Метрики гибкой разработки}

\subsection{Velocity}

\textbf{Velocity} --- количество story points, завершённых командой за спринт.

Если команда завершила за три спринта:
\[
28,\quad 32,\quad 30
\]
story points, то средняя velocity:
\[
V = \frac{28 + 32 + 30}{3} = 30.
\]

Если в бэклоге осталось 120 story points, приблизительное число спринтов:
\[
N = \frac{120}{30} = 4.
\]

Velocity не следует использовать для сравнения разных команд, так как шкалы оценивания локальны.

\subsection{Lead Time и Cycle Time}

\textbf{Lead Time} --- время от появления запроса до его поставки пользователю.

\textbf{Cycle Time} --- время от начала активной работы над задачей до её завершения.

\[
\text{Lead Time} \geq \text{Cycle Time}.
\]

\subsection{Cumulative Flow Diagram}

Кумулятивная диаграмма потока показывает количество задач в разных состояниях во времени. Она помогает обнаруживать узкие места. Если слой \textit{In Progress} растёт, это означает накопление незавершённой работы.

\section{Типичные ошибки внедрения Agile}

\begin{enumerate}
    \item Подмена Agile набором церемоний без изменения культуры.
    \item Отсутствие реального Product Owner.
    \item Слабая инженерная дисциплина: нет тестов, CI, code review.
    \item Использование velocity как инструмента давления.
    \item Слишком большие пользовательские истории.
    \item Отсутствие Definition of Done.
    \item Невозможность команды влиять на способ работы.
    \item Сохранение микроменеджмента под видом Scrum.
    \item Формальные ретроспективы без реальных улучшений.
    \item Игнорирование технического долга.
\end{enumerate}

\section{Agile в современных условиях}

Современные гибкие методологии развиваются под влиянием:
\begin{itemize}
    \item облачной инфраструктуры;
    \item микросервисной архитектуры;
    \item DevOps и платформенной инженерии;
    \item continuous delivery;
    \item продуктового подхода;
    \item data-driven development;
    \item удалённых и распределённых команд;
    \item использования искусственного интеллекта в разработке.
\end{itemize}

При этом базовая идея Agile остаётся прежней: создавать ценность короткими циклами, получать обратную связь и адаптироваться.

\section{Итоговая схема}

\begin{center}
\begin{tikzpicture}[
    node distance=1.5cm,
    box/.style={rectangle, draw, rounded corners, align=center, minimum width=4cm, minimum height=1cm},
    arrow/.style={-{Latex}, thick}
]
\node[box] (values) {Agile-ценности\\и принципы};
\node[box, below left=of values] (scrum) {Scrum\\управление};
\node[box, below=of values] (xp) {XP\\инженерные практики};
\node[box, below right=of values] (kanban) {Kanban\\поток работ};
\node[box, below=of xp] (devops) {DevOps\\поставка и эксплуатация};
\node[box, below=of devops] (value) {Быстрая поставка\\ценности пользователю};

\draw[arrow] (values) -- (scrum);
\draw[arrow] (values) -- (xp);
\draw[arrow] (values) -- (kanban);
\draw[arrow] (scrum) -- (devops);
\draw[arrow] (xp) -- (devops);
\draw[arrow] (kanban) -- (devops);
\draw[arrow] (devops) -- (value);
\end{tikzpicture}
\end{center}

\section{Краткие выводы}

\begin{enumerate}
    \item Agile --- это не конкретный процесс, а система ценностей и принципов, ориентированная на людей, работающий продукт, сотрудничество и адаптацию.
    \item Scrum --- фреймворк эмпирического управления продуктовой разработкой, задающий роли, события и артефакты.
    \item XP --- методология, усиливающая Agile инженерными практиками: TDD, CI, рефакторингом, парным программированием и коллективным владением кодом.
    \item Kanban управляет потоком работ через визуализацию и WIP-лимиты.
    \item Современные гибкие подходы часто комбинируются с DevOps, Lean, continuous delivery и масштабируемыми фреймворками.
    \item Успешное применение Agile требует не только церемоний, но и культуры доверия, технической дисциплины и постоянного улучшения.
\end{enumerate}

\end{document}
